% !TEX root = ../../tree_notes.tex
% =====================================================================
%  2.2  Top-Down Recursion: carrying state down            sec:c02-topdown
%  Exports: def:c02-topdown  lst:c02-topdown-tmpl
%           prob:c02-root-numbers  prob:c02-path-sum-ii  prob:c02-even-grandparent
% =====================================================================
\section{Top-Down Recursion: Carrying State Down}\label{sec:c02-topdown}

\begin{topicintro}
In a \term{top-down} recursion the answer at a node needs information that only its
\emph{ancestors} know -- the depth so far, the running value of the path, a grandparent
pointer -- so we pass that state \emph{in} as a function parameter and refine it on the
way down. This is the mirror image of the height recursion of \secref{c01-anatomy} (which
pulled values \emph{up}); here the values flow from the root to the leaves.

\textbf{Prerequisites.} \secref{c02-model} (the structural-recursion skeleton and the
\nil{} base case) and \defref{c01-rooted} (ancestor, parent, depth).

\textbf{After this section you can}
\begin{itemize}[nosep]
  \item spot a top-down problem: the answer depends on the path \emph{above} the node;
  \item thread an accumulator, a depth, or ancestor pointers through the recursion;
  \item backtrack a shared path buffer correctly (push before, pop after) so siblings do
        not contaminate each other;
  \item solve root-to-leaf numbers, all paths with a target sum, and the
        even-grandparent sum in one clean DFS.
\end{itemize}
\end{topicintro}

\usedin{Chapter 11 (LCA via parent pointers), Chapter 12 (binary lifting carries an
ancestor table down) and Chapter 16 (HLD passes a head node in)}{almost every "path from
the root" or "as a function of depth" query is a top-down DFS; getting the parameter-passing
habit now pays off in all of them.}

\subsection{Intuition}

\begin{analogybox}[Hiking with a running logbook]
Descending a trail from the summit, each camp writes down how many hours it has been since
the top and the list of camps already passed. It does not have to re-count from the bottom:
the number simply arrives from the camp above, gets incremented, and is passed on. That
incremented logbook is the carried state of a top-down DFS.
\end{analogybox}

Consider turning root-to-leaf digits into numbers. When you stand at a node you already know
the value of everything above you (the partial number); you push one more digit down and
add up at the leaves. The information clearly flows downward.

\begin{figure}[H]
\centering
\begin{tikzpicture}[diagram]
  \node[troot] (r) at (0,0) {$4$};
  \node[tnode] (l) at (-1.7,-1.4) {$9$};
  \node[tnode] (c) at (1.7,-1.4) {$0$};
  \node[tleaf] (a) at (-2.5,-2.8) {$5$};
  \node[tleaf] (b) at (-0.9,-2.8) {$1$};
  \draw[tedge] (r) -- (l); \draw[tedge] (r) -- (c);
  \draw[tedge] (l) -- (a); \draw[tedge] (l) -- (b);
  \node[annot, text=primaryblue] at (0,0.9) {$\text{acc}=4$};
  \node[annot, text=primaryblue] at (-1.9,-1.0) {$49$};
  \node[annot, text=primaryblue] at (-2.9,-2.4) {$495$};
  \node[annot, text=primaryblue] at (0.1,-2.4) {$491$};
  \node[annot, text=primaryblue] at (2.2,-1.0) {$40$};
\end{tikzpicture}
\caption{The accumulator $\text{acc}=\text{acc}\times 10+\text{val}$ flows down. Leaves read
$495,491,40$; the answer is their sum $1026$. Nothing flows back up.}\label{fig:c02-topdown}
\end{figure}

\subsection{Definition}

\begin{definition}{Top-down recursion}{c02-topdown}
A recursion $f(\text{node},\ \text{state})$ where \emph{state} is derived from the ancestors
and passed \emph{into} the call. The node computes (or records) its answer from
$\text{state}\cup\{\text{node}\}$, then calls its children with an updated state. It is
correct when the state at a node is exactly the information the path from the root to that
node provides; when state is mutated in place, it must be \term{restored} before the second
child so siblings see the same input.
\end{definition}

\subsection{The Template}

\begin{lstlisting}[caption={General top-down skeleton: refine the carried state, recurse, and undo any in-place change}, label={lst:c02-topdown-tmpl}]
// State is whatever the path above contributes: depth, an accumulator, a vector, ancestors.
void topDown(TreeNode* r, State& s) {
    if (r == nullptr) return;             // base case: the n+1 null links (Lemma 2.1)
    record(s, r);                         // answer this node using the state that arrived
    s.extend(r);                          // refine state for the children (e.g. s.val += ...)
    topDown(r->left,  s);
    topDown(r->right, s);
    s.undo(r);                            // back-track: leave s as you found it for siblings
}
\end{lstlisting}

\vsbox{Top-down vs bottom-up}{Top-down (this section)}{Bottom-up (\secref{c02-bottomup})}{state
flows \emph{in} as a parameter; answer often recorded at leaves or per node}{answer flows
\emph{out} as a return value; combine children to build the parent's answer}

\begin{complexitybox}
A top-down DFS visits all $n$ nodes once -- $\bigO{n}$ -- with $\bigO{h}$ call-stack depth
plus the size of the carried state. If the state is a path \emph{vector copied at every
node}, you pay $\bigO{h}$ per copy and become $\bigO{n h}$; mutate-and-restore keeps it
$\bigO{n}$ with one shared buffer.
\end{complexitybox}

\subsection{Worked Examples}

\begin{example}{Tracing the accumulator on \figref{c02-topdown}}{c02-td-trace-ex}
At the root $\text{acc}=4$. Left child $9$: $\text{acc}=4\cdot 10+9=49$. Its children $5,1$
give $495,491$ (both leaves, so both are added). Right child $0$: $\text{acc}=40$ and it is
itself a leaf. Sum $=495+491+40=1026$.
\end{example}

\begin{tracebox}
For Path Sum II, the shared path buffer evolves as
$[4]\to[4,9]\to[4,9,5]\ \text{(leaf: test)}\ \to[4,9]\to[4,9,1]\to[4]\to[4,0]$ --
the \code{pop\_back} after each child is what stops node $1$ from seeing path $[4,9,5,1]$.
\end{tracebox}

\subsection{Common Pitfalls}

\begin{warningbox}[Forgetting to undo shared state]
If you push onto a shared \code{path} before the left child but forget to pop before
recursing right, the right subtree inherits the left subtree's nodes. Every
mutate-then-recurse needs a matching recurse-then-restore.
\end{warningbox}

\begin{warningbox}[Recomputing what you could carry]
Passing the depth as a parameter is $\bigO{1}$ per node; re-walking from the root to find
the depth at every node is $\bigO{h}$ each and turns an $\bigO{n}$ DFS into $\bigO{n h}$.
If the value is already known on the way in, \emph{carry} it.
\end{warningbox}

\subsection{Solved Problems}

\solvedproblem{c02-root-numbers}{Sum Root to Leaf Numbers}{LeetCode 129}
  {https://leetcode.com/problems/sum-root-to-leaf-numbers/}{LC Medium}
  {top-down accumulator, root-to-leaf}

\begin{problemstatement}[Sum Root to Leaf Numbers]
Each root-to-leaf path spells a decimal number by its node values (e.g.\ path
$4\to9\to5$ is $495$). Return the sum of all such numbers. Leaves are nodes with no
children.
\end{problemstatement}

\solstep{Understanding the Problem}
The number for a path is built digit by digit from the root, so at a node we hold the
partial value $\text{acc}$ and pass $10\cdot\text{acc}+\text{val}$ downward -- the flow of
\figref{c02-topdown}. At a leaf the partial value is a complete number, so it is added to
the total. This is a textbook top-down recursion (\defref{c02-topdown}).

\solstep{Key Observations}
\begin{observation}
$\text{acc}_{\text{child}} = 10\cdot \text{acc}_{\text{parent}} + \text{child->val}$, with
$\text{acc}_{\text{root}}=\text{root->val}$.
\end{observation}
\begin{observation}
Only leaves contribute; an internal node's partial number is never added.
\end{observation}

\solstep{Implementation}
\begin{lstlisting}[caption={Sum Root to Leaf Numbers (LeetCode 129)}, label={lst:c02-root-numbers}]
class Solution {
    int dfs(TreeNode* r, int acc) {
        if (!r) return 0;
        acc = acc * 10 + r->val;                 // carry refined value down
        if (!r->left && !r->right) return acc;   // leaf: complete number
        return dfs(r->left, acc) + dfs(r->right, acc);
    }
public:
    int sumNumbers(TreeNode* root) { return dfs(root, 0); }
};
\end{lstlisting}

\solstep{Complexity}
\begin{complexitybox}
Time $\bigO{n}$, memory $\bigO{h}$. Because the accumulator is passed by \emph{value},
no back-tracking is needed -- each child gets its own copy, which is the cheapest way to
avoid the shared-state pitfall when the state is a single integer.
\end{complexitybox}

\solvedproblem{c02-path-sum-ii}{Path Sum II}{LeetCode 113}
  {https://leetcode.com/problems/path-sum-ii/}{LC Medium}
  {top-down, shared path buffer, backtracking}

\begin{problemstatement}[Path Sum II]
Given a target sum, return every root-to-leaf path whose node values add to it. Each path
is a list of values; a leaf ends a path.
\end{problemstatement}

\solstep{Understanding the Problem}
Unlike \probref{c01-same-tree}, we must \emph{output} the paths, not just decide existence,
so we carry the current path down. Carry the remaining budget too: subtract $\text{val}$ as
we descend and a leaf with remaining $\text{val}$ is a hit. The path is shared state, so it
must be restored (a push/pop DFS exactly as in \figref{c02-topdown}).

\solstep{Key Observations}
\begin{observation}
A leaf is on a valid path iff the target minus all strict ancestors equals its own value.
\end{observation}
\begin{observation}
Record a \emph{copy} of the path at the leaf; the live buffer keeps changing.
\end{observation}

\solstep{Implementation}
\begin{lstlisting}[caption={Path Sum II (LeetCode 113)}, label={lst:c02-path-sum-ii}]
class Solution {
    vector<vector<int>> ans;
    vector<int> path;
    void dfs(TreeNode* r, int rem) {
        if (!r) return;
        path.push_back(r->val);                        // extend state
        rem -= r->val;
        if (!r->left && !r->right && rem == 0) ans.push_back(path);  // leaf hit
        dfs(r->left, rem);
        dfs(r->right, rem);
        path.pop_back();                               // restore for siblings
    }
public:
    vector<vector<int>> pathSum(TreeNode* root, int target) {
        dfs(root, target);
        return ans;
    }
};
\end{lstlisting}

\solstep{Complexity}
\begin{complexitybox}
Time $\bigO{n}$ to visit nodes plus $\bigO{h}$ to copy each found path, worst case
$\bigO{n\cdot h}$ when many paths qualify. Memory $\bigO{h}$ for the live path and stack
(the answer itself can hold $\bigO{n\cdot h}$ values). The \code{pop\_back} is exactly the
"undo shared state" rule -- without it the second child sees the first child's node.
\end{complexitybox}

\solvedproblem{c02-even-grandparent}{Sum of Nodes with Even-Valued Grandparent}{LeetCode 1315}
  {https://leetcode.com/problems/sum-of-nodes-with-even-valued-grandparent/}{LC Medium}
  {top-down, carrying ancestor pointers}

\begin{problemstatement}[Even-Valued Grandparent Sum]
Return the sum of values of nodes whose \emph{grandparent} (two levels up) has an even
value. If no such node exists the sum is $0$.
\end{problemstatement}

\solstep{Understanding the Problem}
The predicate about a node depends on its grandparent -- pure ancestor information, so we
carry grandparent and parent pointers down (\defref{c02-topdown}). At each node, if the
carried grandparent has even value, add the node's value.

\solstep{Implementation}
\begin{lstlisting}[caption={Sum of Nodes with Even-Valued Grandparent (LeetCode 1315)}, label={lst:c02-even-grandparent}]
class Solution {
    int dfs(TreeNode* r, TreeNode* parent, TreeNode* grand) {
        if (!r) return 0;
        int total = (grand && grand->val % 2 == 0) ? r->val : 0;   // decide at this node
        total += dfs(r->left,  r, parent);   // this node becomes parent, old parent becomes grand
        total += dfs(r->right, r, parent);
        return total;
    }
public:
    int sumEvenGrandparent(TreeNode* root) { return dfs(root, nullptr, nullptr); }
};
\end{lstlisting}

\solstep{Complexity}
\begin{complexitybox}
Time $\bigO{n}$, memory $\bigO{h}$. Carrying \code{parent}/\code{grand} as parameters beats
storing them in a map: two pointers per call, updated for free on the way down. (A
bottom-up variant returning grandchild sums also works -- see \secref{c02-bottomup}.)
\end{complexitybox}

\subsection{Practice Problems}

\begin{practiceset}
  \practice{Cousins in Binary Tree}{LeetCode 993}
    {https://leetcode.com/problems/cousins-in-binary-tree/}{LC Easy}
    {carry (depth, parent) down; two nodes are cousins iff equal depth but different parents.}
  \practice{Count Good Nodes in Binary Tree}{LeetCode 1448}
    {https://leetcode.com/problems/count-good-nodes-in-binary-tree/}{LC Medium}
    {a node is good if no ancestor exceeds it; carry the max-so-far down and compare.}
  \practice{Deepest Leaves Sum}{LeetCode 1302}
    {https://leetcode.com/problems/deepest-leaves-sum/}{LC Medium}
    {carry depth down and keep the current deepest level; add values when you reach a new max depth.}
  \practicehh{All Nodes Distance K in Binary Tree}{LeetCode 863}
    {https://leetcode.com/problems/all-nodes-distance-k-in-binary-tree/}{LC Medium}
    {you must go up as well as down -- build parent pointers first, then BFS from the target.}
    {treat the tree as an undirected graph via a parent map; a two-hop BFS reaches both descendants and ancestors within $K$.}
\end{practiceset}
